% TOP_PROBLEM: {"id": 30002802, "title": "Quantum-Group and Affine-Hecke Solutions of Boundary qKZ Equations", "category_id": 16, "category": {"id": 16, "name": "physics", "display_name": "Mathematical Physics", "description": "Problems at the intersection of mathematics and physics.", "slug": "physics", "order_index": 16, "created_at": "2026-07-31T15:26:25.671Z"}, "status": "open", "statusQualification": "This comparison depends on matching parameters, normalizations and analytic domains of the two constructions; an arbitrary basis identification is not the requested relation."}
% FIRST_POSTED: 2026-10-11
% ENTRY_KIND: statement_only
% SUBMISSION: {"kind": "statement_only", "category": "top_problems", "problem_id": "30002802", "proposal_type": "established", "source": "Mini-Workshop: Coideal Subalgebras of Quantum Groups (2015), Oberwolfach Reports, DOI 10.4171/owr/2015/10, pp. 533--569. Primary report accessed 2026-10-11; relevant PDF pages 31. https://doi.org/10.4171/owr/2015/10", "status_evidence": "Canonical source status open retained. Supplied literature triage dated 2026-08-21: The relation between quantum-group correlation or integral solutions and affine-Hecke asymptotic solution bases for qKZ equations still lacks a general theorem. Primary-report access checked 2026-10-11; this historical report and inherited triage do not certify current resolution of every component.", "status_checked_on": "2026-10-11", "duplicate_check": "Compared canonical ID, title and normalized target with the complete local UnsolvedMath registry; normalized exact-text matches: none; title matches: none. Related or overlapping canonical records: none identified in this batch. Primary source context and the existing OWR statement audit used for scope and duplicate review; an exhaustive mathematical-equivalence/literature search was not performed.", "attribution": "Original proposer not named in the supplied canonical record; attribution remains to the identified Oberwolfach contribution/report, not its editors as invented proposers. Explanatory statement preparation and source/semantic audit assisted by OpenAI GPT-6 Codex; no human mathematical authorship or independent proof verification is claimed."}
% SOURCE_STATEMENT: "Determine the relation between the special boundary quantum Knizhnik–Zamolodchikov solutions obtained as quantum correlation functions or $q$-integrals and the affine-Hecke-algebra basis characterized by its asymptotics in a Weyl chamber."
% Preparation and final self-audit: CONTRIBUTING.md prompts C and D; source/audit artifacts retained in the local physics_statements_2026_10_11 batch.
% Imported from: physics_statements.tex; UnsolvedMath ID 30002802
% !TeX program = lualatex
\documentclass[11pt]{article}
\usepackage[margin=25mm]{geometry}
\usepackage{fontspec}
\setmainfont{Latin Modern Roman}
\usepackage{amsmath,amssymb,mathtools}
\usepackage{unicode-math}
\setmathfont{Latin Modern Math}
\usepackage{xurl,hyperref}
\hypersetup{colorlinks=true,urlcolor=blue}
\setcounter{secnumdepth}{0}
\setlength{\parindent}{0pt}
\setlength{\parskip}{5pt}
\setlength{\emergencystretch}{3em}
\newcommand{\sref}[2]{\hyperlink{source:#1:#2}{[S#2]}}
\newcommand{\cataloguescope}[1]{\paragraph{Recorded target.}#1}
\newcommand{\R}{\mathbb R}
\newcommand{\C}{\mathbb C}
\newcommand{\N}{\mathbb N}
\newcommand{\Z}{\mathbb Z}
\newcommand{\Q}{\mathbb Q}
\newcommand{\F}{\mathbb F}
\newcommand{\eps}{\varepsilon}
\begin{document}
% UnsolvedMath ID 30002802; source catalogue code OWR-13495-003
\setcounter{section}{30002801}
\section{Quantum-Group and Affine-Hecke Solutions of Boundary qKZ Equations}\label{30002802}
{\small UnsolvedMath ID 30002802 · Mathematical Physics}
\subsection{Definitions and mathematical statement}

\paragraph{Shared notation.}$\mathbb N=\{1,2,\ldots\}$; the standard number fields have their usual meanings. All additional conventions are stated locally.


\paragraph{Definitions and conventions.}
The boundary quantum Knizhnik--Zamolodchikov (qKZ) equations are compatible difference equations for vector-valued functions of spectral parameters, with exchange operators built from $R$- and boundary $K$-matrices. An affine Hecke algebra is the deformation of an affine Weyl-group algebra subject to braid and quadratic relations. A Weyl chamber is a region determined by inequalities between root coordinates. A $q$-integral is the discrete analogue of integration adapted to multiplicative $q$-shifts.

\paragraph{Target retained from the canonical source catalogue.}
\paragraph{Catalogue wording retained for traceability.}
Determine the relation between the special boundary quantum Knizhnik–Zamolodchikov solutions obtained as quantum correlation functions or $q$-integrals and the affine-Hecke-algebra basis characterized by its asymptotics in a Weyl chamber.

\paragraph{Source scope and qualification.}
This comparison depends on matching parameters, normalizations and analytic domains of the two constructions; an arbitrary basis identification is not the requested relation.
\subsection{Short English statement}
Relate boundary qKZ solutions built by quantum groups or q-integrals to the affine-Hecke basis selected by chamber asymptotics.
\subsection{Sources}
\begin{itemize}
\item[\hypertarget{source:30002802:1}{S1}] ULAM.AI and the UnsolvedMath Contributors, UnsolvedMath, record 30002802 (OWR-13495-003). Catalogue curation under CC BY 4.0; definitions and English exposition added locally.
\url{https://huggingface.co/datasets/ulamai/UnsolvedMath}.
\item[\hypertarget{source:30002802:2}{S2}] Mini-Workshop: Coideal Subalgebras of Quantum Groups (2015), Oberwolfach Reports, DOI 10.4171/owr/2015/10, pp. 533--569. Primary report accessed 2026-10-11; relevant PDF pages 31.
\url{https://doi.org/10.4171/owr/2015/10}.
\end{itemize}
\label{end:30002802}
\end{document}
